% !Mode:: "TeX:UTF-8"

\documentclass[]{WHUBachelor}

\begin{document}

  \miji{}

  \stunum{2015301200666}
  \ctitle{{\LaTeX} Template for Undergraduate Thesis \linebreak Wuhan University \linebreak 毕业论文模板使用指南}
  \cschool{院系名称}
  \cmajor{专业名称}
  \cauthor{作者姓名}
  \cadvisor{指导教师姓名\quad 职称}
  \cdate{二〇一九年五月}

  \CNabstract{
      请使用中文分号“；”分割关键词！\par
      摘要内容摘要内容摘要内容摘要内容摘要内容摘要内容摘要内容摘要内容摘要内容摘要内容摘要内容摘要内容摘要内容摘要内容\par
      摘要内容摘要内容摘要内容摘要内容摘要内容摘要内容摘要内容摘要内容摘要内容摘要内容摘要内容摘要内容摘要内容摘要内容摘要内容摘要内容摘要内容摘要内容摘要内容摘要内容
  }
  \CNkeywords{关键词1；关键词2；关键词3}

  \ENabstract{
      Please use English semicolon and space to separate key words. \par
      This is abstract. This is abstract. This is abstract. This is abstract. This is abstract. This is abstract. This is abstract. This is abstract. \par
      This is abstract. This is abstract. This is abstract. This is abstract. This is abstract. This is abstract. This is abstract. This is abstract. This is abstract. This is abstract. This is abstract. 
  }
  \ENkeywords{key1; key2; key3}

  % !Mode:: "TeX:UTF-8"

%%% ---- 封面 ----- %%%
\setcounter{page}{1}
\begin{titlepage}
  \begin{center}
    \vspace*{10pt}
    {\heiti \zihao{5}%
      \hfill
      \newlength{\myLen}\settowidth{\myLen}{学号\ ：\hspace{3mm}\the\StudentNumber\hspace{3mm}}
      \begin{minipage}[t]{\myLen}
        学号\ ：\uline{\hfill\hspace{3mm}\the\StudentNumber\hspace{3mm}\hfill}  \\[2mm]
        密级\ ：\uline{\hfill\the\miji\hfill}
      \end{minipage}
    }

    \par
    \vspace*{70pt} % 插入空白
    \begin{center}
      \songti \zihao{1} 武汉大学本科毕业论文 \@author
    \end{center}

    \vspace*{54pt}

    % \fbox{
    \begin{minipage}[c][6cm]{12cm}
      \centering
      \heiti \zihao{2} \the\Ctitle
    \end{minipage}

    \par
    \vspace{80pt}
    {\songti\zihao{-3}
      \newcommand\maketabox[1]{\makebox[3.5cm][s]{#1}}
      \begin{tabular}{cp{5cm}}
        \maketabox{院(系)名\ 称\ ：}   & \the\Cschoolname \\[1ex]
        \maketabox{专\ 业\ 名\ 称\ ：} & \the\Cmajor      \\[1ex]
        \maketabox{学\ 生\ 姓\ 名\ ：} & \the\Cauthor     \\[1ex]
        \maketabox{指\ 导\ 教\ 师\ ：} & \the\Csupervisor \\[1ex]
      \end{tabular}
    }
    \par
    \vspace{40pt}
    {
      {\songti \zihao{-2} \the\Cdate}
    }
  \end{center}%
  \clearpage
\end{titlepage}


%%% ---- 郑重申明页 ----- %%%
\newpage
\thispagestyle{empty}
\renewcommand{\baselinestretch}{1.5}  %下文的行距

\vspace*{44pt}
\begin{center}{\ziju{0.5}\songti \zihao{2} \textbf{郑重声明}}\end{center}
\par\vspace*{20pt}
\setlength{\baselineskip}{23pt}
{\zihao{4}
  本人呈交的学位论文，是在导师的指导下，独立进行研究工作所取得的成果，所有数据、图片资料真实可靠。尽我所知，除文中已经注明引用的内容外，本学位论文的研究成果不包含他人享有著作权的内容。对本论文所涉及的研究工作做出贡献的其他个人和集体，均已在文中以明确的方式标明。本学位论文的知识产权归属于培养单位。
  \par
  \vspace*{88pt}
  \hspace*{0.5cm}本人签名: \underline{\hspace{3.5cm}}
  \hspace{2cm}日期: \underline{\hspace{3.5cm}}\hfill\par}


\newcommand\cnkeywords[1]{ {\noindent\heiti\zihao{-4} 关键词: }\zihao{-4}#1}
\newcommand\enkeywords[1]{ {\noindent\bfseries\zihao{-4} Key words: }\zihao{-4}#1}


%%% ---- 中文摘要及关键词 ----- %%%
\newpage
\thispagestyle{empty}
\vspace{10pt}
\begin{center}{\ziju{2}\heiti \zihao{-2} 摘要}\end{center}
\baselineskip=23pt
{\songti \zihao{-4}%

  \the\CNabstract
}
\par
\vspace*{2em}
\cnkeywords{\the\CNkeywords}


%%% ---- 英文摘要及关键词 ----- %%%
\newpage
\thispagestyle{empty}
\vspace{10pt}
\begin{center}{\zihao{-2} \textbf{ABSTRACT}}\end{center}
\baselineskip=23pt
{\zihao{-4}%

\the\ENabstract
}
\par
\vspace*{2em}
\enkeywords{\the\ENkeywords}


  % 正文开始
  \mainmatter

  \chapter{绪论}
  \section{概述}
  与Word等所见即所得编辑工具不同，使用 \LaTeX 工具排版可以将写作与排版过程分离，写作者只需要关心文字的部分，而剩下的排版工作全部交给工具自动完成。
  \section{格式要求}
  正文宋体小4，正文行间距固定为23磅

  通过空一行（两次回车）实现段落换行，仅仅是回车并不会产生新的段落 \par

  也可以通过\verb|\par|命令来新起一段
  \section{各节一级标题}
  我是内容
  \subsection{各节二级标题}
  你是内容
  \subsubsection{各节三级标题}
  他是内容
  \section{字体字号}
  {\songti \bfseries 宋体加粗} {\textbf{English}}

  {\songti \itshape 宋体斜体} {\textit{English}}

  {\songti \bfseries \itshape 宋体粗斜体} {\textbf{\textit{English}}}
  
  \chapter{公式插图表格}
  \section{公式的使用}
  在文中引用公式可以这么写：$a^2+b^2=c^2$这是勾股定理，他还可以表示为$c=\sqrt{a^2+b^2}$，还可以让公式单独一段并且加上编号
  \begin{equation}
  \sin^2{\theta}+\cos^2{\theta}=1 \label{eq:pingfanghe}
  \end{equation}
  还可以通过添加标签在正文中引用公式，如式\eqref{eq:pingfanghe}。

  我们还可以轻松打出一个漂亮的矩阵
  \begin{equation}
    \mathbf{A}=
    \left[\begin{matrix}
      1&2&3&4\\
      11&22&33&44\\
    \end{matrix}\right] \times
    \left[\begin{matrix}
      22&24\\
      32&34\\
      42&44\\
      52&54\\
    \end{matrix}\right]
  \end{equation}
  或者多个带编号的公式

  \begin{eqnarray}
    f_1(x)=12x^2+36x+\sin x\\
    f_2(x)=\sqrt{3}{x^3+3x}
  \end{eqnarray}
  
  \section{插图的使用}
  \LaTeX 环境下可以使用常见的图片格式：JPEG、PNG、PDF、EPS等。当然也可以使用\LaTeX 直接绘制矢量图形。

  图示例如下：
  \begin{figure}[!htb]
    \centering
    \includegraphics[width=0.2\textwidth]
    {figures/whulogo.png}\\
    \caption{插图示例}
    \label{fig:whu}
  \end{figure}

  \verb|[htbp]|选项分别是此处、页顶、页底、独立一页。\verb|[width=\textwidth]|让图片占满整行，或\verb|[width=2cm]|直接设置宽度。可以随时在文中进行引用，如图~\ref{fig:whu}。

  \section{表格的使用}
  表格的输入可能会比较麻烦，可以使用在线的工具，如~\href{https://www.tablesgenerator.com/}{Tables Generator}~能便捷的创建表格。\href{https://en.wikibooks.org/wiki/LaTeX/Tables}{LaTeX/Tables}~上及~\href{https://www.tug.org/pracjourn/2007-1/mori/mori.pdf}{Tables in LaTeX}~也有更多的示例能够参考。

  \subsection{普通表格}
  下面是一些普通表格的示例

  \begin{table}[ht]
    \centering
    \caption{简单表格}
    \label{tab:1}
    \begin{tabular}{|l|c|r|}
      \hline
      我是& 一只 & 普通\\
      \hline
      的& 表格& 呀\\
      \hline
    \end{tabular}
  \end{table}
  
  \begin{table}[ht]
    \centering
    \caption{一般三线表}
    \label{tab:2}
    \begin{tabular}{ccc}
      \hline
      姓名& 学号& 性别\\
      \hline
      张三& 001& 男\\
      李四& 002& 女\\
      \hline
    \end{tabular}
  \end{table}
  
  \subsection{统计表格}
  要创建占满整个文字宽度的表格需要使用到tabularx，如不需要，使用tabular就行。引用表格与其它引用一样，只需要：表~\ref{tab:3}。

  \begin{table}[ht]
    \centering
    \caption{统计数据表格}
    \label{tab:3}
    \begin{tabularx}{\textwidth}{CCCC}
      \toprule
      序号&年龄&身高&体重\\
      \midrule
      1&14&156&42\\
      2&16&158&45\\
      3&14&162&48\\
      4&15&163&50\\
      \cmidrule{2-4} %添加2-4列的中线
      平均&15&159.75&46.25\\
      \bottomrule
    \end{tabularx}
  \end{table}

  \section{列表的使用}
  下面演示了创建有序及无序列表，如需其它样式，\href{https://www.latex-tutorial.com/tutorials/lists/}{LaTeX Lists}~上有更多的示例。

  \subsection{有序列表}
  这是一个计数的列表
	\begin{enumerate}
		\item 第一项
			\begin{enumerate}
				\item 第一项中的第一项
				\item 第一项中的第二项
			\end{enumerate}
		\item 第二项
      \begin{enumerate}[label=(\roman*)]
        \item 第一项中的第一项
        \item 第一项中的第二项
      \end{enumerate}
		\item 第三项
	\end{enumerate}

  \subsection{不计数列表}
	这是一个不计数的列表
	\begin{itemize}
		\item 第一项
		\begin{itemize}
			\item 第一项中的第一项
			\item 第一项中的第二项
		\end{itemize}
		\item 第二项
		\item 第三项
	\end{itemize}

  \chapter{文献引用}

  \section{脚注}
  注释是对论文中特定名词或新名词的注解。注释可用页末注或篇末注的一种。选择页末注的应在注释与正文之间加细线分隔，线宽度为1磅，线的长度不应超过纸张的三分之一宽度。同一页类列出多个注释的，应根据注释的先后顺序编排序号。字体为宋体5号，注释序号以“\circled{1}、\circled{2}”等数字形式标示在被注释词条的右上角。页末或篇末注释条目的序号应按照“\circled{1}、\circled{2}”等数字形式与被注释词条保持一致。示例我需要注释啊\footnote{我是解释注释的}。

  \section{引用文中小节}\label{sec:ref}
  如引用小节~\ref{sec:ref}

  \section{引用参考文献}
  这是一个参考文献引用的范例\cite{kuhn2004man}

  还可以采用上标的引用方式\upcite{江泽民1989能源发展趋势及主要节能措施}

  \chapter{其它格式}
  \section{代码}
  \subsection{原始代码}
  朴实的代码块：

  \begin{verbatim}
    import numpy as np
    print("Hello world!")
  \end{verbatim}

  \subsection{代码风格化}
  代码高亮，更多使用请参考 \href{https://www.overleaf.com/learn/latex/Code_Highlighting_with_minted}{Code Highlighting with minted}

  \usemintedstyle{vs}
  \begin{minted}[linenos,baselinestretch=1.0,frame=lines]{cpp}
    #include <iostream>
    using namespace std;

    int main() 
    {
        cout << "Hello, World!";
        return 0;
    }
  \end{minted}

  \subsection{算法描述/伪代码}
  参考 \href{https://en.wikibooks.org/wiki/LaTeX/Algorithms}{Algorithms}

  % \begin{algorithm}[]
  %   \SetKwInOut{Input}{输入}
  %   \SetKwInOut{Output}{输出}
  %   \caption{\sc DeepWalk\((G, w, d, \gamma, t)\)}\label{DeepWalk}
  %   \Input{图\(G(V,E)\) \\
  %     窗大小\(w\) \\
  %     嵌入维度\(d\) \\
  %     每个节点游走数\(\gamma\) \\
  %     游走距离\(t\)
  %   }
  %   \Output{节点表征矩阵\(\Phi\in\mathbb{R}^{|V|\times d}\)
  %   }
  %   初始化：从\(\mathcal{U}^{|V|\times d}\)采样\(\Phi\) \\
  %   从\(V\)建立二叉树\(T\) \\
  %   \For{\(i=0\) to \(\gamma\)}{
  %     \(\mathcal{O}\) = Shuffle(\(V\)) \\
  %     \ForEach{\(v_i\in\mathcal{O}\)}{
  %       \(\mathcal{W}_{v_i}=RandomWalk(G, v_i, t)\) \\
  %       \(\textsc{SkipGram}(\Phi, \mathcal{W}_{v_i}, w)\)
  %     }
  %   }
  % \end{algorithm}

  
  \section{绘图}
  关于绘图的更多例子，请参考 \href{https://www.overleaf.com/learn/latex/Pgfplots_package}{Pgfplots package}

  % \begin{figure*}[!htb]
  %   \centering
  %   \begin{subfigure}[t]{0.5\textwidth}
  %       \centering
  %       \begin{tikzpicture}
  %         \begin{axis}
  %         \addplot[color=red]{exp(x)};
  %         \end{axis}
  %       \end{tikzpicture}
  %       \caption{2D example}
  %       \label{fig:2d}
  %   \end{subfigure}%
  %   ~ 
  %   \begin{subfigure}[t]{0.5\textwidth}
  %       \centering
  %       \begin{tikzpicture}
  %         \begin{axis}
  %         \addplot3[surf,]
  %         {exp(-x^2-y^2)*x};
  %         \end{axis}
  %       \end{tikzpicture}
  %       \caption{3D example}
  %       \label{fig:3d}
  %   \end{subfigure}
  %   \caption{Figures}
  % \end{figure*}

  % 生成参考文献
  % 使用方法：\bibliography{参考文件1文件名, 参考文献2文件名, ...}
  \bibliography{ref/refs}

  \chapter*{致谢}
  表达对完成论文和学业提供帮助的老师、同学、亲人的感激之情。


  % 下面是附录
  \appendix

  \chapter{这是附录}

  \section{第一个测试}
  测试公式编号
  \begin{equation}
    1+1=2.
  \end{equation}

  表格编号测试

  \begin{table}[h]
    \centering
    \caption{测试表格}
    \begin{tabular}{*{20}c}
      \hline
      11 & 13  & 13  & 13  & 13 \\
      12 & 14  & 13  & 13  & 13 \\
      \hline
    \end{tabular}
  \end{table}

\end{document}